In this part, we study the simplest transducer model, namely Mealy machines. Mealy machines are barely different from deterministic finite automata, and in fact the early literature on automata theory did not even distinguish between the two\footnote{In one of the foundational papers on automata theory, Rabin and Scott comment that they want to simplify previous automata models so that only ``yes'' or ``no'' outputs are possible, see~\cite[p.114]{RabinScott59}.}. From this perspective, Mealy machines are a good starting point for the study of transducers, since many of the results that we care about -- such as continuity, composition and equivalence -- are easy to prove, which we do in Section~\ref{sec:mealy-machines}. There is one exception though, which is the Krohn-Rhodes decomposition theorem from Section~\ref{sec:krohn-rhodes}: of the various decomposition theorems that we will see in this book, this is possibly the most difficult one.